\section{Kalman 滤波：用预测与更新的循环从噪声中估计状态}
\label{sec:16-Control-and-Reinforcement-Learning/01-Optimal-Control/04}

你控制着一间冷库。温度传感器的读数在变——一部分是因为压缩机真的在制冷，一部分是因为有人开门卸货，还有一部分纯粹是传感器的电噪声。你的控制器需要当前的温度状态来做决策，但你拿到的只有被噪声污染的读数。真实的温度躲在传感器后面，你看不到它。前面几节写下的最优控制律，无论是反馈增益还是 HJB 的最优策略，都默认状态可以直接读到；这个默认现在失效了。

这不是一个数据平滑问题。你不能等到收齐了今天全部读数再回头算——控制器需要在每个时刻立刻给出当前的最佳估计，然后马上拿这个估计去算下一拍的控制指令。你要的是一个在线递推：新数据来了，更新一次估计，扔掉旧数据，不必重新遍历历史。

这就是状态估计问题。Kalman filter 用一套简洁的预测-更新循环解决了它，前提是系统是线性的、噪声是 Gaussian 的。但这套前提一旦不满足，filter 给出的估计可能不只是不准，而是系统性地偏到错误的地方去。我们先从它怎么工作讲起，最后再划出它不再可靠的那条线。

\subsection{一个具体的困境：传感器读数不等于状态}

冷库的简化热平衡方程离散化后可以写成：

\[
T_k = a T_{k-1} + b u_k + w_k,
\]

其中 \(T_k\) 是时刻 \(k\) 的仓内温度，\(u_k\) 是压缩机控制量，系数 \(a\in(0,1]\) 描述仓温向环境回归的快慢，\(w_k\) 吸收了墙体蓄热、开门和货物进出中无法用方程精确描述的部分——把它建模为零均值 Gaussian 噪声。

你拿到的传感器读数 \(z_k\) 不等于 \(T_k\)。传感器可能有一个缓慢漂移的零点偏置 \(\beta_k\)，外加短时的随机电噪声 \(\varepsilon_k\)：

\[
z_k = T_k + \beta_k + \varepsilon_k.
\]

如果模型把偏置 \(\beta_k\) 固定为零，filter 会越来越相信自己的温度估计，而对传感器反馈越来越不敏感——最终那个偏置把你带离真实温度，filter 却自信满满地报告一个错误的值。

把偏置也放进状态向量，令 \(x_k = [T_k, \beta_k]^{\trans}\)，并假设偏置按随机游走演化 \(\beta_k = \beta_{k-1} + \zeta_k\)，Kalman filter 就能同时估计温度和偏置。但温度与偏置能否从现有的输入和观测中可靠地分开，取决于模型的结构性质——后文的可观测性——而不是递推步数够不够多。

\subsection{两个高斯，一次联合：更新步骤的心跳}

先写出一般形式的线性 Gaussian 状态空间模型：

\[
X_t = A X_{t-1} + B u_t + W_t,\qquad W_t \sim \mathcal{N}(0, Q),
\]
\[
Y_t = C X_t + V_t,\qquad V_t \sim \mathcal{N}(0, R),
\]

并假设过程噪声、观测噪声与初始状态之间适当独立。

在线性与 Gaussian 的双重假设下，一个关键的闭包性质成立：如果上一时刻的后验是 Gaussian，那么预测分布、联合分布以及更新后的后验全都是 Gaussian。整条递推只需要维护两个量——均值向量和协方差矩阵——不需要追踪整个分布的形状。

假设我们在 \(t-1\) 时刻已经拥有

\[
X_{t-1} \given y_{1:t-1} \sim \mathcal{N}(m_{t-1},\, P_{t-1}).
\]

\textbf{预测步骤}：把上一时刻的后验推过动力学。线性变换加独立 Gaussian 噪声，均值和协方差的传播有闭式：

\[
m_t^- = A m_{t-1} + B u_t,
\]
\[
P_t^- = A P_{t-1} A^{\trans} + Q.
\]

上标 \(-\) 表示``在看见 \(y_t\) 之前''。即使均值按确定性动力学演化，过程噪声 \(Q\) 也会增加协方差——它在说：动力学本身有未建模的部分，时间往前推一步，不确定性的基线就要往上加一截。如果你在实现中忘了加 \(Q\)，filter 会越来越自信，最终拒绝真实观测。

\textbf{更新步骤}：现在 \(y_t\) 到达。你把预测状态和新观测放在一起看它们的联合分布：

\[
\begin{bmatrix}
X_t \\ Y_t
\end{bmatrix}
\sim
\mathcal{N}\!\left(
\begin{bmatrix}
m_t^- \\ C m_t^-
\end{bmatrix},
\begin{bmatrix}
P_t^- & P_t^- C^{\trans} \\
C P_t^- & C P_t^- C^{\trans} + R
\end{bmatrix}
\right).
\]

这个联合协方差矩阵的下三角块 \(C P_t^-\) 是预测状态与预测观测之间的协方差——它度量了``如果真实状态比预测的高一点，观测值也会高多少''。

定义 innovation（新息）：

\[
e_t = y_t - C m_t^-,
\]

即实际观测与预测观测之间的偏差。定义 innovation 协方差：

\[
S_t = C P_t^- C^{\trans} + R.
\]

对联合 Gaussian 应用条件分布公式——给定 \(Y_t = y_t\) 时 \(X_t\) 的条件分布仍然是 Gaussian，其均值和协方差为：

\[
K_t = P_t^- C^{\trans} S_t^{-1},
\]
\[
m_t = m_t^- + K_t e_t,
\]
\[
P_t = P_t^- - K_t C P_t^-.
\]

两个步骤对协方差做的事情正好相反，把它们排在时间轴上，filter 到底在干什么就一目了然。

\begin{center}
\begin{tikzpicture}[>=stealth,line width=0.7pt,x=1.45cm,y=1.25cm]
  % 不确定区间：预测撑开、更新收窄
  \fill[black!8]
    (0,1.225) -- (1,1.414) -- (1,0.816) -- (2,1.080) -- (2,0.734)
    -- (3,1.019) -- (3,0.714) -- (4,1.005) -- (4,0.709) -- (5,1.001)
    -- (5,0.707) -- (6,1.000) -- (6,0.707)
    -- (6,-0.707) -- (6,-1.000) -- (5,-0.707) -- (5,-1.001)
    -- (4,-0.709) -- (4,-1.005) -- (3,-0.714) -- (3,-1.019)
    -- (2,-0.734) -- (2,-1.080) -- (1,-0.816) -- (1,-1.414) -- (0,-1.225)
    -- cycle;
  \draw[black!60]
    (0,1.225) -- (1,1.414) -- (1,0.816) -- (2,1.080) -- (2,0.734)
    -- (3,1.019) -- (3,0.714) -- (4,1.005) -- (4,0.709) -- (5,1.001)
    -- (5,0.707) -- (6,1.000) -- (6,0.707);
  \draw[black!60]
    (0,-1.225) -- (1,-1.414) -- (1,-0.816) -- (2,-1.080) -- (2,-0.734)
    -- (3,-1.019) -- (3,-0.714) -- (4,-1.005) -- (4,-0.709) -- (5,-1.001)
    -- (5,-0.707) -- (6,-1.000) -- (6,-0.707);

  % 估计本身作为基准线
  \draw[dashed,black!55] (-0.1,0) -- (6.3,0);
  \node[anchor=west,font=\footnotesize] at (6.35,0.85) {\(m_t\pm\sqrt{P_t}\)};
  \node[anchor=west,font=\footnotesize] at (6.35,0) {估计 \(m_t\)};

  % 真实状态与估计之差
  \draw[line width=1pt] plot[smooth] coordinates
    {(0,0.45) (0.5,0.85) (1,-0.35) (1.5,-0.60) (2,0.30) (2.5,0.55)
     (3,-0.28) (3.5,-0.50) (4,0.25) (4.5,0.58) (5,-0.30) (5.5,-0.52) (6,0.22)};
  \node[anchor=west,font=\footnotesize] at (6.35,-0.85) {\(x_t-m_t\)};

  % 标注
  \node[anchor=south,font=\footnotesize] at (2.6,1.50) {预测步：加上 \(Q\)，区间撑开};
  \draw[->,black!60] (2.6,1.46) -- (2.6,0.98);
  \node[anchor=north,font=\footnotesize] at (3.0,-1.45) {更新步：吸收 \(y_t\)，区间收窄};
  \draw[->,black!60] (3.0,-1.41) -- (3.0,-1.06);

  % 时间轴
  \draw[->] (-0.15,-2.15) -- (6.5,-2.15);
  \node[anchor=north,font=\small] at (6.5,-2.30) {时间步};
  \foreach \k in {0,1,2,3,4,5,6}{
    \draw (\k,-2.20) -- (\k,-2.10);
    \node[anchor=north,font=\footnotesize] at (\k,-2.22) {\k};
  }
\end{tikzpicture}

\emph{一胀一缩：每一步预测都把不确定区间撑开一点，每一次观测又把它拉回来一截，锯齿最终稳定在一个固定的循环上。真实状态与估计之差落在区间里，是这条带宽有意义的标志。}
\end{center}

锯齿的两条边由两组不同的量控制。上升那一段只与 \(Q\) 有关，是动力学未建模部分持续注入的不确定性；下降那一段的深浅由观测决定，而决定深浅的正是 \(K_t\)。

\(K_t\) 叫 Kalman gain，它是这一整套推导的核心输出。把它拆开来看：\(P_t^- C^{\trans}\) 是状态与观测的协方差——``状态的哪个方向在观测中有投影''；\(S_t^{-1}\) 是观测不确定性的逆——``观测有多可信''。二者相乘的结果决定了：在观测告诉你的方向上，你用多大的步长修正预测，也就是上图里那道下降沿有多陡。

标量情形把这件事写得最干净。此时 \(K_t = P_t^-/(P_t^- + R)\)，正是预测方差在新息总方差中占的份额，它只依赖 \(P_t^-\) 与 \(R\) 的比值。传感器很吵时 \(R\) 大，份额小，新读数被打折，估计基本沿着动力学走；预测本身很不确定时份额接近 1，一次观测就能把估计拉过去一大截。极端情形也是连续的：\(R\) 趋于零而预测不确定性不为零时，\(K_t \to C^{-1}\)（\(C\) 可逆的情形），估计几乎完全跟随观测，动力学只负责在两次观测之间填空。多维情形无非是把``比值''换成了矩阵意义下的相对精度。

\begin{center}
\begin{tikzpicture}[>=stealth,line width=0.8pt,x=1cm,y=1cm]
  % 三条不确定区间
  \draw[black!70] (0.3,2.4) -- (2.7,2.4);
  \draw[black!70] (0.3,2.22) -- (0.3,2.58);
  \draw[black!70] (2.7,2.22) -- (2.7,2.58);
  \fill (1.5,2.4) circle (2.1pt);

  \draw[black!70] (4.9,1.5) -- (6.1,1.5);
  \draw[black!70] (4.9,1.32) -- (4.9,1.68);
  \draw[black!70] (6.1,1.32) -- (6.1,1.68);
  \fill (5.5,1.5) circle (2.1pt);

  \draw[line width=1.1pt] (4.16,0.6) -- (5.24,0.6);
  \draw[line width=1.1pt] (4.16,0.42) -- (4.16,0.78);
  \draw[line width=1.1pt] (5.24,0.42) -- (5.24,0.78);
  \fill (4.7,0.6) circle (2.1pt);

  \node[anchor=east,font=\small] at (-0.15,2.4) {预测};
  \node[anchor=east,font=\small] at (-0.15,1.5) {观测};
  \node[anchor=east,font=\small] at (-0.15,0.6) {融合};
  \node[anchor=west,font=\footnotesize] at (6.3,2.4) {半宽 \(\sqrt{P_t^-}\)};
  \node[anchor=west,font=\footnotesize] at (6.3,1.5) {半宽 \(\sqrt{R}\)};
  \node[anchor=west,font=\footnotesize] at (6.3,0.6) {半宽 \(\sqrt{P_t}\)};

  % 垂线落到数轴
  \draw[dashed,black!45] (1.5,2.4) -- (1.5,-0.5);
  \draw[dashed,black!45] (5.5,1.5) -- (5.5,-0.5);
  \draw[dashed,black!45] (4.7,0.6) -- (4.7,-0.5);

  % 两段份额
  \draw[<->,black!65] (1.5,0.0) -- (4.7,0.0);
  \draw[<->,black!65] (4.7,0.0) -- (5.5,0.0);
  \node[anchor=south,font=\footnotesize] at (3.1,0.04) {\(K_t\)};
  \node[anchor=west,font=\footnotesize] at (5.65,0.0) {\(1-K_t\)};

  % 数轴
  \draw[->] (-0.3,-0.5) -- (7.6,-0.5);
  \node[anchor=north,font=\small] at (7.6,-0.62) {状态取值};
  \foreach \p in {1.5,4.7,5.5}{\draw (\p,-0.58) -- (\p,-0.42);}
  \node[anchor=north,font=\footnotesize] at (1.5,-0.62) {\(m_t^-\)};
  \node[anchor=north,font=\footnotesize] at (4.72,-0.62) {\(m_t\)};
  \node[anchor=north,font=\footnotesize] at (5.62,-0.9) {\(y_t\)};
\end{tikzpicture}

\emph{每条区间的半宽就是它自己的不确定性。融合后的估计落在预测与观测之间、偏向窄的那一侧，而且比两者都窄。它落在什么位置由 \(K_t\) 决定，\(K_t\) 又只是两个宽度的比——增益不是调出来的参数，是被 \(P_t^-\) 与 \(R\) 算出来的。}
\end{center}

Kalman gain 不是在两个固定权重之间做选择——它会随每一步的 \(P_t^-\) 动态变化。filter 刚启动时，初始协方差很大，gain 很大，快速吸收早期数据；稳定运行后，\(P_t\) 收敛到稳态值，gain 也趋于常数，锯齿图上表现为最右边那几个周期高度不再变化。

同一套递推还有一条纯概率的进路：把整条状态序列与观测序列看成一个大的联合 Gaussian，链式的条件独立结构保证条件化每次只牵涉相邻两片，于是预测与更新就是 Gaussian 条件公式沿时间轴的反复调用。那条路线在 \hyperref[sec:13-Machine-Learning/07-Graphical-and-Sequential-Models/03]{``Kalman filter 是动态系统上的 Gaussian 条件化''} 里走完；它从分布的结构进入，本节从控制器的需求进入，两个角度落在同一组公式上。

\subsection{走一个小规模数值例子}

设一维位置跟踪问题。状态只有位置 \(x_t\)，动力学是随机游走：\(x_t = x_{t-1} + w_t\)，\(w_t \sim \mathcal{N}(0, 0.1)\)。观测直接看位置加噪声：\(y_t = x_t + v_t\)，\(v_t \sim \mathcal{N}(0, 1)\)。

初始估计：\(m_0 = 0\)，\(P_0 = 10\)（对初始位置高度不确定）。

第一步预测：\(m_1^- = 0\)，\(P_1^- = 10 + 0.1 = 10.1\)。第一次观测为 \(y_1 = 1.2\)，于是

\[
K_1 = \frac{10.1}{10.1 + 1} \approx 0.91,
\qquad
m_1 = 0 + 0.91 \times 1.2 \approx 1.09,
\qquad
P_1 = 10.1 - 0.91 \times 10.1 \approx 0.91.
\]

协方差从 10.1 骤降到 0.91——第一次观测带来了大量信息。gain 0.91 说明 filter 几乎完全信任这次观测（观测噪声 1 相对于预测不确定性 10.1 确实很小）。

再做一步。预测：\(m_2^- = 1.09\)，\(P_2^- = 0.91 + 0.1 = 1.01\)。假设 \(y_2 = 1.5\)。gain：\(K_2 = 1.01 / (1.01 + 1) \approx 0.50\)。更新：\(m_2 \approx 1.30\)，\(P_2 \approx 0.50\)。gain 降为 0.50——经过第一步，filter 已经比较确定了，对后续观测的采纳更谨慎。这正是锯齿图最左边那两个周期：第一次下降沿极深，之后迅速变浅。

几步之后 \(P_t\) 趋于稳态。把 \(A = C = 1\) 代进代数 Riccati 方程

\[
P = A P A^{\trans} + Q - A P C^{\trans}(C P C^{\trans} + R)^{-1} C P A^{\trans}
\]

得到 \(P^2 - 0.1P - 0.1 = 0\)，正根 \(P \approx 0.370\)，这是稳态的预测协方差；对应的后验协方差与稳态增益都等于 \(0.370/1.370 \approx 0.270\)。无论跑多少步，协方差不会再低于它，因为每一步的过程噪声都在持续注入新的不确定性——锯齿的上升沿永远不会消失。

这个例子暴露了 Kalman filter 的权衡本质：\(Q\) 和 \(R\) 的相对大小决定了 filter 最终在信任动力学与信任观测之间停在哪里。如果 \(Q\) 设得太小，filter 对动力学过于自信，观测修正的效果被压制，估计会滞后于真实状态的变化——这称为估计滞后（estimator lag），是手动调低 \(Q\) 最常见的副作用。

\subsection{可观测性：有些方向再多数据也看不见}

Kalman filter 的协方差 \(P_t\) 会缩小，但不是所有方向都缩得一样快。某个状态方向如果长期不通过 \(C X_t\) 影响观测，filter 就收不到关于它的信息。动力学 \(A\) 可以把该方向的能量逐步耦合到可观测方向——但也可能永远不耦合，让那个方向一直藏在暗处。

设状态维数为 \(d\)，线性系统的 observability 矩阵

\[
\mathcal{O} =
\begin{bmatrix}
C \\ CA \\ \vdots \\ CA^{d-1}
\end{bmatrix}
\]

回答了这个问题：如果 \(\mathcal{O}\) 满列秩，在理想无噪条件下，有限时间段内的观测足以区分任意两个不同的初始状态。如果不满秩，则存在不可观测的子空间——其中的状态分量无论采集多少数据都无法从观测中恢复，即使 \(R\) 任意小。

冷库的例子把这条判据算得出来。状态是温度与偏置，\(A = \diag(a, 1)\)，读数是两者之和，\(C = [1\ \ 1]\)，于是

\[
\mathcal{O} = \begin{bmatrix} 1 & 1 \\ a & 1 \end{bmatrix},
\qquad
\det \mathcal{O} = 1 - a .
\]

判据落在 \(a\) 上。若 \(a = 1\)，即仓温本身也是随机游走、没有任何向环境回归的趋势，行列式为零：温度与偏置只有和可测、差不可测，采多久的数据都分不开这两者。若 \(a < 1\)，温度会衰减而偏置不会，两者的时间特征不同，系统在理论上可观测——但 \(a\) 越接近 1，\(\mathcal{O}\) 越接近奇异，把两者分开所需的信噪比越高，有限时间窗内的偏置估计仍可能毫无意义。不是递推算错了，是信息在数据里太弱。

对偶的概念是可控性——通过控制输入 \(u_t\) 能否把状态驱动到任意目标，取决于 controllability 矩阵的秩。可控性和可观测性一起构成了线性系统理论的骨架：可控性保证你能把系统带到想去的地方，可观测性保证你知道系统当前在什么地方。

\subsection{Filtering、prediction、smoothing 不是同一件事}

同一个模型可以问三种问题，区别在于允许动用的数据窗口。Filtering 算 \(p(x_t \given y_{1:t})\)，只用当前和过去的观测，这是在线控制唯一合法的选择。Prediction 算 \(p(x_{t+h} \given y_{1:t})\)，从当前往后推 \(h\) 步，中间没有新观测进来，协方差只会按预测步的规律一路撑大。Smoothing 算 \(p(x_t \given y_{1:T})\) 且 \(T > t\)，等整段数据到齐后回头修正过去的状态估计，它的方差严格不大于滤波。

离线轨迹重建可以放心用 smoothing——你已经看到全程数据了，允许未来观测修正你对过去的判断。但实时控制不能在时刻 \(t\) 使用尚未发生的 \(y_{t+1:T}\)。如果评价模型时误把平滑后的状态当成在线特征来用，就会出现时间泄漏：模型在训练时偷看了未来，上线后做不到同样的精度，而离线指标漂亮得让人不忍心怀疑。

\subsection{数值实现不是套公式就行了}

写在纸上的公式用 \(S_t^{-1}\) 是为了表达线性算子的逆作用。真实的数值代码应该解线性系统

\[
S_t Z = C P_t^-
\]

而不是显式求逆——求逆比求解慢且数值不稳定。

协方差更新的简式 \(P_t = P_t^- - K_t C P_t^-\) 在有限精度运算下可能丢失对称性或半正定性。更稳健的 Joseph form 同时从左右两侧应用修正：

\[
P_t = (I - K_t C) P_t^- (I - K_t C)^{\trans} + K_t R K_t^{\trans},
\]

右端是两个半正定项之和，对称与非负是写在形式里的，不靠运算精度去维持。对于病态问题，平方根 filter 直接传播 Cholesky 因子而非协方差矩阵本身，进一步避免舍入误差导致的负``方差''。数值稳定性不是模型之外的附属问题：负方差或非对称协方差会污染下一次更新中的概率计算，让 filter 在没有任何模型错误的情况下自我毁坏。

还有一种信息形式的 Kalman filter，传播的不是 \((m_t, P_t)\)，而是信息矩阵 \(\Lambda_t = P_t^{-1}\) 与信息向量 \(\eta_t = P_t^{-1} m_t\)。它在初始不确定性极大的场景里有天然优势：\(P_0 \to \infty\) 对应 \(\Lambda_0 \to 0\)，用零信息矩阵表达``一无所知''比用一个巨大的协方差数值安全得多。在传感器网络和多机器人协作估计中，多个独立观测对信息矩阵和信息向量的贡献是可加的，分布式融合因此变得便捷。

\subsection{非线性与非 Gaussian：闭式递推在此终结}

Kalman filter 的优雅来自一个很窄的前提：线性变换把 Gaussian 映成 Gaussian，联合 Gaussian 的条件分布仍是 Gaussian。这两个性质缺一个，闭式递推就会断裂。

如果动力学或观测是非线性的，Gaussian 经过非线性变换后通常不再是 Gaussian。Extended Kalman filter 的做法是在当前均值处取一阶 Taylor 展开——用局部的线性 \(A_t = \partial f / \partial x\) 和 \(C_t = \partial h / \partial x\) 替代固定的 \(A\) 和 \(C\)。近似的质量取决于：在当前不确定性椭球范围内，真正的非线性函数偏离其线性近似有多远。不确定性很大时，一阶近似可能完全不可靠。一个经典的 EKF 失效场景是：机器人初始位姿不确定性很大，而观测模型（比如距离和方位的极坐标转换）在远离线性化点处高度非线性——此时 EKF 的线性化可能指向完全错误的更新方向。

Unscented Kalman filter 不取导数，而是用一组确定性选择的 sigma points 穿过非线性函数，再用变换后的样本点重建均值和协方差——它相当于在矩的层面做了至少二阶精度的近似。Particle filter 则完全放弃 Gaussian 假设，用带权的离散粒子表示任意形状的后验分布。

这三种方法之间的层级很清楚：Kalman 假设线性加 Gaussian，得到闭式解；EKF 与 UKF 放开线性，退到近似的矩匹配；Particle filter 放开 Gaussian，改用 Monte Carlo 表示任意分布。选择哪一种，取决于你的模型在这条层层放宽的阶梯上站在哪一级——而不是取决于哪一个的代码看起来更现代。

多峰后验是一个特别值得警惕的信号。如果状态确实有两种以上截然不同的可能取值（比如机器人可能在走廊的这一侧或那一侧，但绝不可能在墙里面），用单峰 Gaussian 去近似会把概率质量平均到不可能的区域。这时候即使线性化再精确，被表示族的结构限制也会让 filter 给出一个物理上不可能的估计。问题是表示能力的不足，不是参数精度的不足。Particle filter 也不能解决所有问题——粒子在高维状态空间中的退化（绝大部分粒子的权重趋于零）是它自己的核心困难，需要靠重采样和好 proposal 分布来缓解。

Kalman filter 的力量来源于一组非常具体且严苛的闭包条件：线性映射、Gaussian 噪声与 Markov 状态。看清楚了这些条件，也就看清楚了它在什么情况下是可靠的工具，在什么情况下必须换模型而不只是调协方差参数。

回到开头那间冷库。现在你有了当前状态的最佳估计，但控制器要的是控制量，把估计直接当成真状态塞进反馈增益，凭什么是对的？这一步的合法性不能靠直觉，得靠一条定理。\hyperref[sec:16-Control-and-Reinforcement-Learning/01-Optimal-Control/05]{``LQG：看不全状态时怎样继续反馈''} 会证明，在同一组线性 Gaussian 假设下，估计与反馈可以分开设计而不损失最优性——本节的锯齿与前面几节的最优增益，恰好在那里接成一个完整的闭环。
